%======================================================================
\chapter{Correspondence with the formal proof}\label{app:formal}
%======================================================================

The table lists the definitions and results of this volume that carry a ``Formal counterpart''
line, together with the corresponding declarations of the Lean development. All declaration
names start with \lean{HPC.}. The formal statements are sometimes slightly more general, or are
split differently into lemmas; the table points to the place where the content of each result is
proved. Rows are in the order of the text.

The main theorems of the volume are, in the Lean development (prefix \lean{HPC.} omitted):
\begin{center}
\begin{tabular}{@{}p{0.42\textwidth}p{0.54\textwidth}@{}}
\toprule
Characterization of the arc (Corollary~\ref{cor:arc-char}) & \lean{Moore.isIrreducibleConnected_iff_arc}\\
Characterization of the circle (Theorem~\ref{thm:circle-char}) & \lean{Moore.circle_char}\\
Quotient of a Moore decomposition is metrizable (Theorem~\ref{thm:quot-metrizable}) & \lean{Moore.quotient_metrizable_plane}\\
Mullikin's theorem, corrected (Theorem~\ref{thm:mullikin}) & \lean{Moore.mullikin_theorem5}\\
Moore 1925, Theorem~11 (Theorem~\ref{thm:m25-11}) & \lean{Moore.m25_theorem11}\\
Jordan curve theorem (Theorem~\ref{thm:jct}) & \lean{Jones.jordan_curve_theorem}\\
Moore 1919, Theorem~D (Theorem~\ref{thm:theoremD}) & \lean{Moore.m19_theoremD}\\
A compact grid is a square (Theorem~\ref{thm:grid-square}) & \lean{Sq.Grid.exists_square}\\
Schoenflies theorem in a Moore space (Theorem~\ref{thm:schoenflies}) & \lean{Moore.m19_schoenflies_boundary_proof}\\
Moore's characterization of the plane (Theorem~\ref{thm:moore-plane}) & \lean{Moore.m19_plane}\\
Moore 1925, Theorem~21 (Theorem~\ref{thm:m25-21}) & \lean{Moore.m25_theorem21}\\
Moore 1925, Theorem~22 (Theorem~\ref{thm:m25-22}) & \lean{Moore.m25_theorem22}, \lean{Moore.m25_theorem22_moore}\\
Monotone images (Corollary~\ref{cor:monotone-plane}) & \lean{Hagopian.exists_plane_quotient_of_monotone}\\
Schoenflies theorem in the plane (Theorem~\ref{thm:plane-schoenflies}) & \lean{Hagopian.plane_schoenflies_aux}\\
Two-sided collars (Theorem~\ref{thm:collar}) & \lean{Hagopian.exists_collar}\\
\bottomrule
\end{tabular}
\end{center}
All of them are proved in the Lean development without \lean{sorry} and without axioms beyond the
standard ones of Lean and Mathlib. The printed forms that are false (Section~\ref{sec:intro-diff}) are
refuted by the declarations whose names begin with \lean{not_}.

\begin{longtable}{@{}p{0.40\textwidth}p{0.56\textwidth}@{}}
\toprule
In these notes & In the Lean development\\
\midrule
\endhead
\bottomrule
\endfoot
\multicolumn{2}{@{}l}{\emph{Chapter~\ref{ch:arcs}}}\\
Definition~\ref{def:arc} & \lean{HPC.Moore.Gen.TArc}, \lean{HPC.Moore.Gen.TSCC}.\\
Sub-arcs (Lemma~\ref{lem:subarc}) & \lean{HPC.Moore.Gen.TArc.image_Icc}, \lean{HPC.Moore.Gen.TArc.exists_sub}, \lean{HPC.Moore.Gen.TArc.nonempty_diff}, \lean{HPC.Moore.Gen.TArc.isPreconnected_diff}, \lean{HPC.Moore.Gen.TArc.subset_closure_diff}.\\
Lemma~\ref{lem:conn-frontier} & \lean{HPC.Moore.Gen.conn_sub_of_frontier}, \lean{HPC.Moore.Gen.conn_sub_ext_of_frontier}.\\
First points, gluing, joining (Lemma~\ref{lem:arc-ops}) & \lean{HPC.Moore.Gen.TArc.first_point}, \lean{HPC.Moore.Gen.TArc.union}, \lean{HPC.Moore.Gen.TArc.join}.\\
Lemma~\ref{lem:insert-conn} & \lean{HPC.Moore.ac_insert_isPreconnected}.\\
Cuts (Lemma~\ref{lem:cut}) & \lean{HPC.Moore.ac_compl_isPreconnected}, \lean{HPC.Moore.ac_cut}.\\
Proposition~\ref{prop:arc-order} & \lean{HPC.Moore.acL}, \lean{HPC.Moore.acL_asymm}, \lean{HPC.Moore.acL_trans}, \lean{HPC.Moore.acL_total}, \lean{HPC.Moore.acL_ray}, \lean{HPC.Moore.ac_exists_between}.\\
Irreducible continua are arcs (Theorem~\ref{thm:irr-arc}) & \lean{HPC.Moore.ac_exists_strictMono}, \lean{HPC.Moore.ac_homeomorph}.\\
Arc characterization (Corollary~\ref{cor:arc-char}) & \lean{HPC.Moore.isIrreducibleConnected_iff_arc}.\\
Definition~\ref{def:cutarc} & \lean{HPC.Moore.IsCutArc}.\\
Proposition~\ref{prop:cutarc} & (a) \lean{HPC.Moore.IsCutArc.irreducible}; (b) \lean{HPC.Moore.IsCutArc.exists_embedding}, \lean{HPC.Moore.isCutArc_range}, \lean{HPC.Moore.isCutArc_image_Icc}; (c) \lean{HPC.Moore.eq_image_Icc_of_subset_range}, \lean{HPC.Moore.ends_of_image_Icc}; (d) \lean{HPC.Moore.exists_nonCut_diff}.\\
Moore 1920, Theorem~4 (Theorem~\ref{thm:circle-char}) & \lean{HPC.Moore.circle_noCut}, \lean{HPC.Moore.circle_half}, \lean{HPC.Moore.circle_char}.\\
Proposition~\ref{prop:scc-cut} & \lean{HPC.Moore.Gen.tscc_cut}; \lean{HPC.Jones.glue_circle}, \lean{HPC.Jones.tscc_homeomorph_circle}.\\
Limit arcs (Theorem~\ref{thm:limit-arc}) & \lean{HPC.Moore.Gen.tarc_of_seq}; with trimming of arcs that meet more than consecutively, \lean{HPC.Moore.Gen.exists_tarc_of_seq}.\\
Replacement of a null family of sub-arcs (Theorem~\ref{thm:arc-replace}) & \lean{HPC.Moore.Gen.arc_replace}.\\[4pt]
\multicolumn{2}{@{}l}{\emph{Chapter~\ref{ch:usc}}}\\
Definition~\ref{def:moore-decomp} & \lean{HPC.Moore.MooreDecomp}, \lean{HPC.Moore.elt}, \lean{HPC.Moore.ptsOf}.\\
Definition~\ref{def:tube} & \lean{HPC.Moore.tube}.\\
Proposition~\ref{prop:closed-map} & \lean{HPC.Moore.isClosedMap_mk_plane}, \lean{HPC.Moore.isOpen_tube}, \lean{HPC.Moore.mem_nhds_iff_tube}, \lean{HPC.Moore.usc_of_isClosedMap}.\\
Theorem~\ref{thm:quot-metrizable} & \lean{HPC.Moore.t1Space_quot}, \lean{HPC.Moore.regularSpace_quot}, \lean{HPC.Moore.secondCountable_quot}, \lean{HPC.Moore.quotient_metrizable_plane}.\\
Proposition~\ref{prop:pts-conn} & \lean{HPC.Moore.isPreconnected_ptsOf}, \lean{HPC.Moore.ptsOf_connectedComponentIn}, \lean{HPC.Moore.isOpen_connectedComponentIn_quot}, \lean{HPC.Moore.locallyConnectedSpace_quot}, \lean{HPC.Moore.frontier_eq_image}.\\
Corollary~\ref{cor:quot-tools} & \lean{HPC.Moore.quot_brouwer}, \lean{HPC.Moore.quot_janiszewski}, \lean{HPC.Moore.m25_unbounded_unique}, \lean{HPC.Moore.isConnected_compl_singleton_quot}.\\
Proposition~\ref{prop:moore-notions} & \lean{HPC.Moore.isLimitElt_iff}, \lean{HPC.Moore.m25_domain_iff}, \lean{HPC.Moore.m25_boundary_iff}.\\
Moore, Theorem~1 (Theorem~\ref{thm:m25-1}) & \lean{HPC.Moore.m25_theorem1}.\\
Definition~\ref{def:moore-region} & \lean{HPC.Moore.mooreRegion}, \lean{HPC.Moore.spaceOfElts}.\\
Moore, Theorem~3 (Theorem~\ref{thm:m25-3}) & \lean{HPC.Moore.exists_region_in_tube}, \lean{HPC.Moore.m25_theorem3}.\\
Moore, Theorem~4 (Theorem~\ref{thm:m25-4}) & \lean{HPC.Moore.isTopologicalBasis_mooreRegion}, \lean{HPC.Moore.m25_theorem4}, \lean{HPC.Moore.spaceOfElts_topology}, \lean{HPC.Moore.cl_eq_closure_quot}.\\
Moore, Theorem~2 (Theorem~\ref{thm:m25-2}) & \lean{HPC.Moore.m25_theorem2}.\\
Proposition~\ref{prop:p420} & \lean{HPC.Moore.m25_axioms1245}, \lean{HPC.Moore.m25_FA_1_15}.\\[4pt]
\multicolumn{2}{@{}l}{\emph{Chapter~\ref{ch:elements}}}\\
Definition~\ref{def:arc-elt} & \lean{HPC.Moore.IsArcElt}, \lean{HPC.Moore.IsSCCElt}.\\
Proposition~\ref{prop:arc-elt-top} & \lean{HPC.Moore.isCompact_of_isClosed_bounded}, \lean{HPC.Moore.isArcElt_iff_tarc}, \lean{HPC.Moore.isSCCElt_iff_tscc}.\\
Moore, Theorems 5--8 and 10 (Theorem~\ref{thm:m25-5-10}) & \lean{HPC.Moore.m25_theorem5}, \lean{HPC.Moore.m25_theorem6}, \lean{HPC.Moore.m25_theorem7}, \lean{HPC.Moore.m25_theorem8}, \lean{HPC.Moore.m25_theorem10}.\\
Moore, Theorem~9 (Theorem~\ref{thm:m25-9}) & \lean{HPC.Moore.m25_theorem9}.\\
Corollary~\ref{cor:pts-nonsep} & \lean{HPC.Moore.not_separatesPlane_ptsOf}.\\
Moore, Theorem~12 (Theorem~\ref{thm:m25-12}) & \lean{HPC.Moore.m25_theorem12}.\\
Text of Section~\ref{sec:points} & \lean{HPC.Moore.mooreDecomp_ptSetoid}, \lean{HPC.Moore.ptHomeo}.\\[4pt]
\multicolumn{2}{@{}l}{\emph{Chapter~\ref{ch:jordan}}}\\
At least two components (Proposition~\ref{prop:mullikin-two}) & \lean{HPC.Moore.not_isConnected_compl_union}.\\
At most two components; Eilenberg (Proposition~\ref{prop:mullikin-three}) & \lean{HPC.Moore.exists_jump}, \lean{HPC.Moore.no_log_pow_of_bounded_component}, \lean{HPC.Moore.not_two_bounded_components}, \lean{HPC.Moore.compl_union_two_components}.\\
Mullikin's theorem (Theorem~\ref{thm:mullikin}) & \lean{HPC.Moore.mullikin_theorem5}.\\
Lemma~\ref{lem:frontier-scc} & \lean{HPC.Moore.frontier_eq_of_scc}.\\
Moore, Theorem~11 (Theorem~\ref{thm:m25-11}) & \lean{HPC.Moore.m25_theorem11}, \lean{HPC.Moore.sccInterior_spec}.\\
Jordan curve theorem (Theorem~\ref{thm:jct}) & \lean{HPC.Jones.jordan_curve_theorem}.\\[4pt]
\multicolumn{2}{@{}l}{\emph{Chapter~\ref{ch:axioms}}}\\
Moore's derived notions (Definition~\ref{def:moore-notions}) & \lean{HPC.Moore.Space}, \lean{HPC.Moore.Space.IsLimitPt}, \lean{HPC.Moore.Space.bdry}, \lean{HPC.Moore.Space.cl}, \lean{HPC.Moore.Space.ext}, \lean{HPC.Moore.Space.IsConn}, \lean{HPC.Moore.Space.IsBoundedM}, \lean{HPC.Moore.Space.IsDomain}, \lean{HPC.Moore.Space.IsArc}, \lean{HPC.Moore.Space.IsSCC}, \lean{HPC.Moore.Space.IsSimpleChain}.\\
Moore's axioms, F.A.\ pp.~132--133 (Definition~\ref{def:moore-axioms}) & \lean{HPC.Moore.Space.Axiom1} \dots\ \lean{HPC.Moore.Space.Axiom8}, \leansub{HPC.Moore.Space.Axiom5}{1}, \leansub{HPC.Moore.axiom5}{1}\lean{_of_axiom5}.\\
The bridge (Proposition~\ref{prop:bridge}) & \lean{HPC.Moore.Space.topology}, \lean{HPC.Moore.fa_theorem6}, \lean{HPC.Moore.isTopologicalBasis_region}, \lean{HPC.Moore.isLimitPt_iff}, \lean{HPC.Moore.cl_eq_closure}, \lean{HPC.Moore.isClosedM_iff}, \lean{HPC.Moore.isConn_iff}, \lean{HPC.Moore.isDomain_iff}, \lean{HPC.Moore.heineBorel_iff}, \lean{HPC.Moore.t1Space}, \lean{HPC.Moore.regularSpace}, \lean{HPC.Moore.fa_theorem4}, \lean{HPC.Moore.secondCountable}, \lean{HPC.Moore.metrizable}, \lean{HPC.Moore.isCompact_closure_region}, \lean{HPC.Moore.isBoundedM_of_isCompact}.\\
F.A.\ Theorems~1--3, 5, 7--9, 11--14, 17--19, 22 (Theorem~\ref{thm:fa-elementary}) & \lean{HPC.Moore.fa_theorem1}, \lean{HPC.Moore.fa_theorem2}, \lean{HPC.Moore.fa_theorem3}, \lean{HPC.Moore.fa_theorem5}, \lean{HPC.Moore.fa_theorem7}, \lean{HPC.Moore.fa_theorem8}, \lean{HPC.Moore.fa_theorem9}, \lean{HPC.Moore.fa_theorem11}, \lean{HPC.Moore.fa_theorem12}, \lean{HPC.Moore.fa_theorem13}, \lean{HPC.Moore.fa_theorem14}, \lean{HPC.Moore.fa_theorem17}, \lean{HPC.Moore.fa_theorem18}, \lean{HPC.Moore.fa_theorem19}, \lean{HPC.Moore.isConn_univ}, \lean{HPC.Moore.fa_theorem22}.\\
F.A.\ Theorem~10 (Theorem~\ref{thm:fa10}) & \lean{HPC.Moore.fa_theorem10}, \lean{HPC.Moore.exists_simpleChain_sublist}.\\
F.A.\ Theorem~15 (Theorem~\ref{thm:fa15}) & \lean{HPC.Moore.exists_chains}, \lean{HPC.Moore.exists_irreducible}, \lean{HPC.Moore.fa_theorem15}.\\
F.A.\ Theorem~16 (Corollary~\ref{cor:fa16}) & \lean{HPC.Moore.fa_theorem16}.\\
Proposition~\ref{prop:moore-arcs} & \lean{HPC.Moore.isArc_iff}, \lean{HPC.Moore.isSCC_iff_tscc}.\\[4pt]
\multicolumn{2}{@{}l}{\emph{Chapter~\ref{ch:curves}}}\\
F.A.\ Theorems~20 and~21 (Theorem~\ref{thm:fa20-21}) & \lean{HPC.Moore.fa_theorem20}, \lean{HPC.Moore.fa_theorem21}.\\
The interior of a closed curve (Corollary~\ref{cor:interior}) & \lean{HPC.Moore.existsUnique_interior}, \lean{HPC.Moore.jordan}, \lean{HPC.Moore.Space.IsInteriorOf}.\\
Text of the section ``The interior of a closed curve'' & \lean{HPC.Moore.Gen.ArcConn}, \lean{HPC.Moore.Gen.Th21Prop}, \lean{HPC.Moore.Gen.Ax8Prop}, \lean{HPC.Moore.Gen.TInt}.\\
Text of the section ``The interior of a closed curve'' & \lean{HPC.Moore.Gen.exists_access}, \lean{HPC.Moore.Gen.exists_cross}.\\
F.A.\ Theorem~23 (Theorem~\ref{thm:fa23}) & \lean{HPC.Moore.fa_theorem23}.\\
F.A.\ Theorem~24 (Theorem~\ref{thm:fa24}) & \lean{HPC.Moore.fa_theorem24}, \lean{HPC.Moore.Gen.th24_sub}, \lean{HPC.Moore.Gen.th24_ext}, \lean{HPC.Moore.Gen.th24_disj}.\\
F.A.\ Theorem~25 (Theorem~\ref{thm:fa25}) & \lean{HPC.Moore.fa_theorem25}, \lean{HPC.Moore.Gen.th25}.\\
F.A.\ Theorem~26 (Theorem~\ref{thm:fa26}) & \lean{HPC.Moore.fa_theorem26}, \lean{HPC.Moore.Gen.th26}.\\
F.A.\ Theorem~27 (Theorem~\ref{thm:fa27}) & \lean{HPC.Moore.fa_theorem27}, \lean{HPC.Moore.Gen.th27}.\\
Splitting an interior (Corollary~\ref{cor:split-interior}) & \lean{HPC.Moore.split_interior}, \lean{HPC.Moore.split_interior_of}.\\
F.A.\ Theorem~28 (Theorem~\ref{thm:fa28}) & \lean{HPC.Moore.fa_theorem28}.\\
Text of the section ``Arcs near a curve: Theorems 28--30'' & \lean{HPC.Moore.fa_theorem29}, \lean{HPC.Moore.fa_theorem30}.\\
F.A.\ Theorem~31; the pocket construction (Theorem~\ref{thm:fa31}) & \lean{HPC.Moore.fa_theorem31}, \lean{HPC.Moore.th31_pocket}.\\
The two sides (Corollary~\ref{cor:sides}) & \lean{HPC.Moore.sides}.\\
Propagation and accessibility; F.A.\ Theorem~32 (Theorem~\ref{thm:fa32}) & \lean{HPC.Moore.side_propagate}, \lean{HPC.Moore.accessible}, \lean{HPC.Moore.fa_theorem32}.\\
F.A.\ Theorems~33 and~34 (Theorem~\ref{thm:fa33-34}) & \lean{HPC.Moore.reach_end}, \lean{HPC.Moore.joinZ}, \lean{HPC.Moore.fa_theorem33}, \lean{HPC.Moore.fa_theorem34_region}, \lean{HPC.Moore.fa_theorem34}.\\
F.A.\ Theorems~39 and~40; accessibility of a closed curve (Theorem~\ref{thm:fa39-40}) & \lean{HPC.Moore.access}, \lean{HPC.Moore.fa_theorem39}, \lean{HPC.Moore.fa_theorem40}.\\
F.A.\ Theorems~35 and~36 (Theorem~\ref{thm:fa35-36}) & \lean{HPC.Moore.fa_theorem35}, \lean{HPC.Moore.fa_theorem36}, \lean{HPC.Moore.fa_theorem36_interior_subset}.\\
Chords form a null family; Moore 1919, Theorem~A (Lemma~\ref{lem:chords-null}) & \lean{HPC.Moore.chords_null}, \lean{HPC.Moore.m19_theoremA}.\\
Text of the section ``Merging interiors: Theorems 37--43'' & \lean{HPC.Moore.exists_maximal_chord}, \lean{HPC.Moore.fa_theorem37}, \lean{HPC.Moore.fa_theorem38}.\\
F.A.\ Theorem~41 (Theorem~\ref{thm:fa41}) & \lean{HPC.Moore.fa_theorem41_caseI_nested}, \lean{HPC.Moore.fa_theorem41_caseI}, \lean{HPC.Moore.fa_theorem41_of_subset}, \lean{HPC.Moore.fa_theorem41_caseII}, \lean{HPC.Moore.Gen.arc_replace}, \lean{HPC.Moore.fa_theorem41}.\\
F.A.\ Theorem~42 (Theorem~\ref{thm:fa42}) & \lean{HPC.Moore.fa_theorem42_of_41}, \lean{HPC.Moore.fa_theorem42}.\\
F.A.\ Theorem~43, corrected (Theorem~\ref{thm:fa43}) & \lean{HPC.Moore.dual_component}, \lean{HPC.Moore.fa_theorem43}.\\[4pt]
\multicolumn{2}{@{}l}{\emph{Chapter~\ref{ch:rulings}}}\\
Definition~\ref{def:quad} & \lean{HPC.Moore.Space.Quad}.\\
Definition~\ref{def:double-ruling} & \lean{HPC.Moore.Space.DoubleRuling}, \lean{HPC.Moore.Space.DoubleRuling.carrier}, \lean{HPC.Moore.Space.DoubleRuling.cells}.\\
Text of the section ``Quadrilaterals and double rulings'' & \lean{HPC.Moore.Space.ArcRuling}, \lean{HPC.Moore.Space.ArcRuling.toDoubleRuling}, \lean{HPC.Moore.Space.DoubleRuling.toArcRuling}, \lean{HPC.Moore.Space.ArcRuling.param_mono}.\\
Lemma~\ref{lem:tiling} & \lean{HPC.Moore.Space.Tiling}, \lean{HPC.Moore.Space.Tiling.add_arc}, \lean{HPC.Moore.Space.Tiling.add_arc_finite}, \lean{HPC.Moore.Space.Tiling.components}.\\
Subdivision lemma (Proposition~\ref{prop:cells-count}) & \lean{HPC.Moore.Space.DoubleRuling.tiling}, \lean{HPC.Moore.doubleRuling_cells_ncard}.\\
Text of the section ``Cutting, gluing, extending and refining'' & \lean{HPC.Moore.Space.Quad.splitX}, \lean{HPC.Moore.Space.Quad.rot}, \lean{HPC.Moore.Space.Quad.Cut}.\\
Cut and glue (Lemma~\ref{lem:cut-glue}) & \lean{HPC.Moore.Space.Quad.Cut.restrictA}, \lean{HPC.Moore.Space.Quad.Cut.restrictB}, \lean{HPC.Moore.Space.Quad.Cut.glue}, \lean{HPC.Moore.Space.ArcRuling.extendY}, \lean{HPC.Moore.Space.ArcRuling.extendY_finite}.\\
Text of the section ``Cutting, gluing, extending and refining'' & \lean{HPC.Moore.Space.Rulable}, \lean{HPC.Moore.Space.DoubleRulable}.\\
Refinement (Proposition~\ref{prop:refine}) & \lean{HPC.Moore.Space.ArcRuling.refine}, \lean{HPC.Moore.doubleRuling_refine}, \lean{HPC.Moore.Space.Quad.Cut.rulable}.\\
Definition~\ref{def:chord-quad} & \lean{HPC.Moore.Space.IsChord}, \lean{HPC.Moore.Space.Crosses}, \lean{HPC.Moore.Space.IsCrossCut}.\\
Text of the section ``Theorem C'' & \lean{HPC.Moore.no_ruling_of_two_branches}, \lean{HPC.Moore.tangent_configuration}, \lean{HPC.Moore.not_m19_theoremC_printed}.\\
Chords (Proposition~\ref{prop:chords-rulable}) & \lean{HPC.Moore.Space.Pocket}, \lean{HPC.Moore.chordWeight}, \lean{HPC.Moore.outer_tile}, \lean{HPC.Moore.exists_gap}, \lean{HPC.Moore.chords_rulable}.\\
A curve within (Proposition~\ref{prop:inside-rulable}) & \lean{HPC.Moore.inside_rulable}.\\
Text of the section ``Theorem C'' & \lean{HPC.Moore.extract_chords}, \lean{HPC.Moore.curve_rulable}, \lean{HPC.Moore.m19_theoremC_corrected}, \lean{HPC.Moore.m19_theoremC_crosscuts}, \lean{HPC.Moore.m19_theoremC_within}.\\
A nice curve about a compact set (Proposition~\ref{prop:nice-curve}) & \lean{HPC.Moore.lens_data}, \lean{HPC.Moore.lens_family}, \lean{HPC.Moore.exists_nice_curve}.\\
Moore 1919, Theorem~D (Theorem~\ref{thm:theoremD}) & \lean{HPC.Moore.rulable_of_cover}, \lean{HPC.Moore.m19_theoremD}.\\
Theorem~D for regions (Corollary~\ref{cor:theoremD-regions}) & \lean{HPC.Moore.m19_theoremD_regions_of_curves}, \lean{HPC.Moore.m19_theoremD_regions}.\\[4pt]
\multicolumn{2}{@{}l}{\emph{Chapter~\ref{ch:plane}}}\\
Definition~\ref{def:family-lines} & \lean{HPC.Sq.Family}, \lean{HPC.Sq.PreGrid}, \lean{HPC.Sq.box}, \lean{HPC.Sq.Grid}.\\
A compact grid is a square (Theorem~\ref{thm:grid-square}) & \lean{HPC.Sq.Family.pos}, \lean{HPC.Sq.Grid.dense_t}, \lean{HPC.Sq.Grid.mem_L_of_pos_eq}, \lean{HPC.Sq.Grid.continuousOn_pos}, \lean{HPC.Sq.Grid.eq_of_pos_eq}, \lean{HPC.Sq.Grid.image_pos}, \lean{HPC.Sq.Grid.exists_square}, \lean{HPC.Sq.Grid.exists_square_param}.\\
Boxes are components (Lemma~\ref{lem:boxes}) & \lean{HPC.Sq.PreGrid.box_nonempty}, \lean{HPC.Sq.PreGrid.box_sub_component}.\\
Nested rulings (Lemma~\ref{lem:nested-rulings}) & \lean{HPC.Moore.nested_rulings}, \lean{HPC.Moore.nested_rulings_fund}.\\
Schoenflies; Moore 1919, Theorem~E (Theorem~\ref{thm:schoenflies}) & \lean{HPC.Moore.lineFam}, \lean{HPC.Moore.lr_mono}, \lean{HPC.Moore.preGrid}, \lean{HPC.Moore.preGrid_box_sub_cell}, \lean{HPC.Moore.m19_schoenflies_proof}, \lean{HPC.Moore.m19_schoenflies_param}, \lean{HPC.Moore.m19_schoenflies_boundary_proof}, \lean{HPC.Moore.m19_schoenflies_boundary}.\\
Corollary~\ref{cor:schoenflies-curve} & \lean{HPC.Moore.exists_quad}, \lean{HPC.Moore.exists_square_of_interior}.\\
Annulus lemma (Theorem~\ref{thm:annulus}) & \lean{HPC.Moore.exists_cuts}, \lean{HPC.Moore.annulus_decomp}, \lean{HPC.Moore.exists_collar}, \lean{HPC.Moore.m19_annulus}, \lean{HPC.Moore.m19_annulus_square}.\\
Lemma~\ref{lem:curve-about-compact} & \lean{HPC.Moore.preconnectedSpace}, \lean{HPC.Moore.exists_curve_about_compact}.\\
Moore 1919, Theorem~F (Theorem~\ref{thm:theoremF}) & \lean{HPC.Moore.m19_theoremF_proof}, \lean{HPC.Moore.m19_theoremF}.\\
Pasting (Proposition~\ref{prop:paste}) & \lean{HPC.Sq.exhaust_homeomorph}, \lean{HPC.Sq.plane_of_collars}.\\
Moore's characterization of the plane (Theorem~\ref{thm:moore-plane}) & \lean{HPC.Moore.m19_plane_proof}, \lean{HPC.Moore.m19_plane}.\\
Moore 1919, Theorem~G and the final paragraph (Corollary~\ref{cor:theoremG}) & \lean{HPC.Moore.m19_theoremG}, \lean{HPC.Moore.m19_openCurve}, \lean{HPC.Moore.m19_final}.\\[4pt]
\multicolumn{2}{@{}l}{\emph{Chapter~\ref{ch:quotient}}}\\
Interiors and exteriors (chapter introduction) & \lean{HPC.Moore.sccInterior}, \lean{HPC.Moore.sccExterior}, \lean{HPC.Moore.sccInterior_spec}, \lean{HPC.Moore.closure_sccInterior}, \lean{HPC.Moore.closure_sccExterior}.\\
Proposition~\ref{prop:elt-top} & \lean{HPC.Moore.isArcElt_iff_tarc}, \lean{HPC.Moore.isSCCElt_iff_tscc}, \lean{HPC.Moore.isSCCElt_union}, \lean{HPC.Moore.t2Space_quot}, \lean{HPC.Moore.arcConn_quot}, \lean{HPC.Jones.tscc_homeomorph_circle}.\\
Definition~\ref{def:restricted} & \lean{HPC.Moore.restrictedRegion}, \lean{HPC.Moore.spaceRestricted}.\\
Lemma~\ref{lem:restricted-moore} & \lean{HPC.Moore.mooreRegion_of_restricted}, \lean{HPC.Moore.th21Prop_restricted}, \lean{HPC.Moore.ax8Prop_restricted}, \lean{HPC.Moore.isCompact_closure_restricted}, \lean{HPC.Moore.not_isCompact_univ_quot}.\\
Janiszewski's theorem for elements (Lemma~\ref{lem:quot-jan}) & \lean{HPC.Moore.quot_janiszewski}.\\
Moore, Theorem~13 (Theorem~\ref{thm:m25-13}) & \lean{HPC.Moore.not_subset_arc_of_cut}, \lean{HPC.Moore.m25_theorem13}.\\
Moore, Theorem~14 (Theorem~\ref{thm:m25-14}) & \lean{HPC.Moore.exists_tscc_of_arcs}, \lean{HPC.Moore.m25_theorem14}.\\
Moore, Theorems~15 and~16 (Theorem~\ref{thm:m25-15-16}) & \lean{HPC.Moore.m25_theorem15}, \lean{HPC.Moore.m25_theorem16}.\\
Moore, Theorem~17 (Theorem~\ref{thm:m25-17}) & \lean{HPC.Moore.m25_theorem17}.\\
Moore, Theorem~18 (Theorem~\ref{thm:m25-18}) & \lean{HPC.Moore.m25_theorem18}.\\
Moore, Theorem~19 (Theorem~\ref{thm:m25-19}) & \lean{HPC.Moore.m25_theorem19}.\\
Moore, Theorem~20 (Theorem~\ref{thm:m25-20}) & \lean{HPC.Moore.isTopologicalBasis_restricted}, \lean{HPC.Moore.m25_theorem20}, \lean{HPC.Moore.spaceRestricted_topology}.\\
Moore, Theorem~21 (Theorem~\ref{thm:m25-21}) & \lean{HPC.Moore.m25r_axiom1}, \lean{HPC.Moore.m25r_axiom2}, \lean{HPC.Moore.m25r_axiom3}, \lean{HPC.Moore.m25r_axiom4}, \lean{HPC.Moore.m25r_axiom5}, \lean{HPC.Moore.m25r_axiom6}, \lean{HPC.Moore.m25r_axiom7}, \lean{HPC.Moore.m25r_axiom8}, \lean{HPC.Moore.m25_axioms6'7'}, \lean{HPC.Moore.m25_theorem21}.\\
Moore, Theorem~22 (Theorem~\ref{thm:m25-22}) & \lean{HPC.Moore.m25_theorem22}, \lean{HPC.Moore.m25_theorem22_moore}.\\
Proposition~\ref{prop:monotone} & \lean{HPC.Hagopian.isPreconnected_preimage_of_monotone}, \lean{HPC.Hagopian.isDecomposableSpace_of_monotone_image}.\\
Lemma~\ref{lem:fibre-usc} & \lean{HPC.Hagopian.fibreSetoid}, \lean{HPC.Hagopian.fibre_usc}.\\
Corollary~\ref{cor:monotone-plane} & \lean{HPC.Hagopian.fibreSetoid_mooreDecomp}, \lean{HPC.Hagopian.exists_plane_quotient_of_monotone_aux}, \lean{HPC.Hagopian.exists_plane_quotient_of_monotone}.\\
Cross-cuts in the plane (Theorem~\ref{thm:plane-split}) & \lean{HPC.Hagopian.plane_split}.\\
Schoenflies theorem in the plane, one side prescribed (Theorem~\ref{thm:plane-schoenflies}) & \lean{HPC.Hagopian.plane_schoenflies_aux}.\\
Two-sided collars (Theorem~\ref{thm:collar}) & \lean{HPC.Hagopian.IsCollar}, \lean{HPC.Hagopian.exists_collar}.\\

\end{longtable}
